\section{Case Studies and Proof-of-Concept Synergies}\label{sec:case_studies}

The case studies in this section are intended to demonstrate that the
TCDA tuple of \Cref{def:tcda} is not an empty abstraction: in each of
five concrete settings we can write down the four components
$(\Mc, G, \Tc, \Cc)$ explicitly, locate the example in the taxonomy of
\Cref{fig:taxonomy}, and indicate which parts are already established in
the literature versus which constitute the gaps formalised in
\Cref{sec:research}. The five cases were chosen to (i)~cover all four
top-level domains A--D of the taxonomy, (ii)~engage each of the four
seed works of the field---\citet{ibeling2021topological,
kim2026topological,lecci2014statistical,elyaagoubi2023statistical}---in
at least one case study, and (iii)~span what we expect to be the most
active sites of near-term TCDA research: biomedical imaging and
molecular structure (\Cref{sec:cs_tate}), multivariate brain time series
(\Cref{sec:cs_cycles}), latent-variable confounding diagnostics
(\Cref{sec:cs_confounding}), Mapper-based causal abstraction with
calibrated uncertainty (\Cref{sec:cs_mapper}), and spatiotemporal
climate attribution (\Cref{sec:cs_climate}). \Cref{tab:cs_summary}
summarises the mapping back to the taxonomy.

\subsection{Case Study 1: Topological average treatment effects on
biomedical imaging and molecular structure}\label{sec:cs_tate}

\paragraph{Setting.}
Modern biomedical imaging and computational chemistry deliver
non-Euclidean, high-dimensional outcome data---cubical CT volumes,
volumetric MR scans, three-dimensional molecular point clouds---on
which the classical ATE of \cref{eq:ate} is ill-posed because the
outcome is not a scalar but a shape \citep{soares2020sars,
axelrod2022geom,iqbal2025classification,carlsson2020topological}. The
treatment effect of interest is a contrast between two shape
distributions, not between two means. The first complete TCDA solution
to this problem is the topological average treatment effect of
\citet{kim2026topological}.

\paragraph{TCDA instantiation.}
We instantiate \Cref{def:tcda} as follows. (i)~$\Mc$ is the space of
imaging or molecular outcomes, equipped with a natural ambient metric:
the Hausdorff or Wasserstein distance on $\R^3$ point clouds for
molecular data, the $L^\infty$ pixel distance for cubical images.
(ii)~$G$ is a small ``treatment-to-outcome'' DAG with a binary
treatment $A \in \{0,1\}$ (e.g.\ COVID-positive versus
COVID-negative; ligand-bound versus apo conformation), pre-treatment
covariates $X$ (age, sex, scanner protocol, molecular descriptors),
and the shape-valued outcome $Y \in \Mc$. (iii)~$\Tc$ is composition of
two stable functors: persistent homology
$\PHk{k}$---cubical for images \citep{kaczynski2004computational},
Vietoris--Rips for molecular point clouds---followed by the power-
weighted silhouette $\varphi(\cdot; D, r)$ of \cref{eq:silhouette}.
By \Cref{lem:silhouette_stability}, $\Tc$ is Lipschitz from
$(\mathcal{P}(\Mc), W_1)$ into $L^\infty(T)$. (iv)~$\Cc$ is the
\emph{topological average treatment effect}
\begin{equation}
  \TATE \;=\; \E\bigl[\varphi(\cdot;\,\Tc(Y^1),r)\bigr]
              \;-\; \E\bigl[\varphi(\cdot;\,\Tc(Y^0),r)\bigr],
  \label{eq:tate_silhouette}
\end{equation}
together with the Wasserstein contrast
$W_1(\Dgmk{k}(P^1),\Dgmk{k}(P^0))$ on the underlying diagrams.

\paragraph{Estimation and inference.}
Under the standard identifying assumptions (C1)--(C3) of
\Cref{sec:prelim_ci}, \citet{kim2026topological} construct a doubly-
robust augmented inverse-probability-weighting estimator for
\cref{eq:tate_silhouette} in which the outcome nuisance is a
function-on-vector regression of $\varphi(\cdot; \Tc(Y), r)$ on
$(X,A)$ and the propensity nuisance is $\widehat\pi(x) = \widehat\Prob(A=1\mid X = x)$.
Under product-rate nuisance conditions the estimator achieves a
$\sqrt n$ functional rate in $L^\infty(T)$ and admits a tight
weak-convergence limit, from which a level-$\alpha$ test of
$H_0 : W_1(\Dgmk{k}(P^1),\Dgmk{k}(P^0)) = 0$ is read off. The
statistical scaffolding inherits directly from
\citet{lecci2014statistical}'s programme of confidence sets and
bootstrap procedures for persistence diagrams, landscapes, and
silhouettes; without this scaffolding, the TCDA estimator would deliver
a point estimate but no uncertainty.

\paragraph{Stability.}
The estimator inherits the quantitative robustness of persistent
homology via \Cref{prop:stability-transfer}: a Hausdorff perturbation
of the support of $P^a$ translates into an $L^\infty$ perturbation of
the estimated silhouette TATE, with an explicit constant tracking the
silhouette weight $r$. This is the concrete payoff of Domain B1:
where classical AIPW estimators are sensitive to imaging artefacts and
scanner drift, the topological AIPW estimator is provably robust to
small Hausdorff perturbations of the support.

\paragraph{Why this case study matters and what remains open.}
The case study completely populates one ``corner'' of the TCDA
taxonomy---sub-domains A3 and B1, with applications in D3---and is the
clearest existence proof that the framework is non-vacuous: a
non-Euclidean outcome that previously admitted only descriptive TDA
now admits a calibrated, doubly-robust causal estimator. Remaining
gaps targeted in \Cref{sec:research} include the scaling of
$\PHk{k}$ to volumetric three-dimensional inputs, transportability of
the estimator across imaging protocols (a positivity issue), and the
extension of the construction to non-binary or continuous treatments.

\subsection{Case Study 2: Cycle-based causal discovery in
multivariate dependence networks}\label{sec:cs_cycles}

\paragraph{Setting.}
Multivariate time-series data---neural recordings, financial returns,
climate indices---admit a natural \emph{dependence network} whose edges
are weighted by pairwise statistical dependence (mutual information,
partial coherence, transfer entropy). Causal questions on these
networks concern the presence and direction of feedback loops:
re-entrant cortical loops in EEG \citep{sizemore2018cliques,
elyaagoubi2023statistical}, contagion cycles in financial markets
\citep{gidea2018topological,ismail2022early}, oceanic--atmospheric
teleconnections in climate \citep{strommen2025topological}. Persistent
homology of the dependence network in homology degree~$1$ is the
natural detector of such loops.

\paragraph{TCDA instantiation.}
We set (i)~$\Mc$ = the space of multivariate windowed time series in
$\R^d$ with the supremum-distance topology; (ii)~$G$ = a directed
Granger or transfer-entropy graph encoding lagged causal influence
among the $d$ variables; (iii)~$\Tc$ = the composition ``windowed
dependence network'' $\to$ ``filtration by edge weight'' $\to$
$\PHk{1}$, optionally smoothed by the silhouette of
\cref{eq:silhouette}; (iv)~$\Cc$ = a hypothesis test on the number of
persistent $1$-cycles in $\Dgmk{1}$, or---more ambitiously---a recovery
test for the cycle structure of $G$ from $\Dgmk{1}$. The construction
is exactly the one analysed in \citet{elyaagoubi2023statistical}, whose
mixture of latent AR(2) processes allows researchers to \emph{prescribe}
the number of cycles in the network and so to evaluate $\Tc$-based
discovery procedures against ground truth.

\paragraph{Inference and benchmark.}
\citet{elyaagoubi2023statistical} construct simulation-based confidence
regions for $\PHk{1}$ of the dependence network and use them to test
hypotheses about cycle count. The companion work
\citet{elyaagoubi2023topological} extends the construction to
asymmetric dependence and so to genuinely directed Granger graphs. The
case study realises sub-domains A1 (PH for causal structure learning),
C1 (Betti numbers as causal explanation), and D1 (time-series
causality) of \Cref{sec:framework} in a single, fully executable
pipeline. \citet{bando2022causal,harnack2017causal,gholizadeh2018short,
krasnovsky2025causal} pursue closely related strategies for dynamical
causal links via PH of Takens embeddings, providing a complementary
source of empirical evidence.

\paragraph{From a cycle to a causal mechanism.}
Sub-domain C1 raises the harder question: when does a persistent
$1$-cycle deserve the label ``causal feedback loop'' rather than the
weaker label ``cyclic correlation pattern''? The dependence-network
construction does not, on its own, supply directionality. A natural
answer, partial but concrete, is provided by \Cref{prop:identifiability}
specialised to the dependence-network functor $\Tc$: if the assignment
$G \mapsto \Tc(\mathcal{P}_G)$ is injective on a class of admissible
graphs (e.g.\ acyclic Granger graphs with non-Gaussian innovations,
in which case classical identifiability via the LiNGAM family applies
\citep{shimizu2006lingam}), then $\PHk{1}$-based discovery is
identifying. Where this injectivity condition fails---e.g.\ for two
distinct cyclic graphs that yield the same windowed dependence
network---the $1$-cycle is a \emph{joint} causal pattern of an
equivalence class of graphs, and downstream identification requires
auxiliary assumptions. The simulation framework of
\citet{elyaagoubi2023statistical} makes both successful and unsuccessful
cases controllable in the laboratory, and we list the
characterisation of injectivity as an explicit open avenue in
\Cref{sec:research}.

\paragraph{Why this case study matters and what remains open.}
The case study shows that an existing TDA benchmark is essentially a
TCDA benchmark in waiting: equipping the dependence-network filtration
with an explicit causal graph $G$ converts the topological cycle
detector into a topological causal-discovery procedure, with calibrated
uncertainty inherited from \citet{lecci2014statistical}'s programme.
The remaining methodological gap is the formal identifiability theory
that would convert the construction from a benchmark-validated heuristic
into a procedure with deductive guarantees on real (non-simulated)
multivariate time series.

\subsection{Case Study 3: Confounding diagnostics via persistent
residual topology}\label{sec:cs_confounding}

\paragraph{Setting.}
Suppose an analyst posits a causal graph $G$ over an observed set of
variables and fits the corresponding SCM. Under correct specification
the residuals are independent of the predicted endogenous variables
and---crucially for TCDA---their joint support carries no topological
information beyond Euclidean noise. When an unobserved confounder
$U$ acts on two or more observed variables along an unobserved orbit,
the residual support inherits a persistent low-dimensional topology
that no orientation of edges within $G$ can explain. This observation
is the basis of \emph{topological confounding diagnostics}
(sub-domain B2 of \Cref{sec:framework}) and connects, via
\citet{ibeling2021topological}, to the topological no-free-lunch
theorem of the field.

\paragraph{TCDA instantiation.}
(i)~$\Mc$ = the residual space of a fitted SCM, viewed as a point cloud
in $\R^{|V|}$; (ii)~$G$ = the analyst's hypothesised graph, treated as
a null hypothesis; (iii)~$\Tc$ = $\PHk{1}$ of the Vietoris--Rips
filtration of the residual cloud, with bootstrap confidence sets
\citep{fasy2014confidence,chazal2013bootstrap}; (iv)~$\Cc$ = a test of
the null hypothesis ``the persistent $1$-classes of $\Tc$ have
lifetimes consistent with white-noise residuals.''

\paragraph{Worked example: the circular confounder.}
Let $U$ be uniformly distributed on the unit circle $S^1$ and define
$X = \cos U + \epsilon_X$ and $Y = \sin U + \epsilon_Y$ with
$\epsilon_X, \epsilon_Y$ small, independent, mean-zero noise. An
analyst who assumes no edge $X \to Y$ and no confounder will correctly
estimate $\widehat{\ATE}(X \to Y) \approx 0$, since on average $X$
carries no information about $Y$. The same analyst, however, will
discover a persistent $1$-cycle in $\Dgmk{1}$ of the $(X,Y)$ residuals
whose lifetime exceeds any plausible noise band, signalling the action
of $U$. Formally:

\begin{proposition}[Topological confounding signature]
\label{prop:confounding_signature}
Let $U$ be a continuous random variable supported on a topological
space $\mathcal{U}$ with non-trivial $H_1(\mathcal{U};\R) \neq 0$, and
let $X_v = f_v(U,\epsilon_v)$ for $v \in V$ with $f_v$ continuous and
$\epsilon_v$ Gaussian and independent of $U$. Suppose the analyst's
hypothesised graph $G$ omits $U$ and is otherwise correct. Then the
limit support of the joint residuals admits a continuous surjection
onto $\mathcal{U}$, and the persistent $\PHk{1}$ of any sufficiently
fine Vietoris--Rips filtration of the residuals contains at least one
class whose expected lifetime is bounded away from zero uniformly in
the noise scale.
\end{proposition}

\begin{proof}[Sketch]
By the continuous-mapping theorem and the assumed continuity of each
$f_v$, the image of $U$ under $v \mapsto X_v$ is a continuous
surjection $\mathcal{U} \to \mathrm{supp}(P^{X})$. The corresponding
homology map $H_1(\mathcal{U}) \to H_1(\mathrm{supp}(P^X))$ is
non-trivial whenever the $(f_v)$ collectively are not homotopic to a
constant, which is generic. The stability theorem
(\cref{eq:stability}) ensures that the corresponding cycle persists in
the empirical $\PHk{1}$ of finite samples, with lifetime bounded below
by the homological diameter of $\mathcal{U}$. Full details, including
the rate at which the lifetime concentrates, follow the stability
arguments of \citet{cohen2007stability,chazal2016structure}.
\end{proof}

\paragraph{Discussion: detection is not identification.}
\Cref{prop:confounding_signature} is one-sided: a persistent residual
$1$-cycle is consistent with confounding but does not, on its own,
recover $U$ or even identify the variables it acts on. This asymmetry
is exactly the content of the topological causal hierarchy theorem of
\citet{ibeling2021topological}: the set of SCMs admitting
assumption-free identification of $P(y\mid \doop(x))$ is topologically
meagre, so almost any SCM carries residual topological structure that
\emph{flags} confounding but cannot \emph{certify} its absence. The
sheaf-theoretic programme of \citet{mahadevan2021causal,
mahadevan2022universal,mahadevan2025sheaf} provides a complementary
reading in which the residual cycle is interpreted as a non-trivial
class in a sheaf-cohomology obstruction to identification; this is
the natural home of sub-domain B3.

\paragraph{Why this case study matters and what remains open.}
The case study converts an existing TDA descriptor---persistent
$1$-cycles of a point cloud---into a diagnostic that is logically
linked to a causal property (residual confounding) via
\Cref{prop:confounding_signature}, and shows that the limits of the
diagnostic are precisely those identified in
\citet{ibeling2021topological}'s meagre-set theorem. Operationalising
the test on real data---calibrating the lifetime threshold against
noise, controlling for nuisance topology induced by sampling, and
extending to higher-dimensional confounders with $H_k(\mathcal{U}) \neq 0$
for $k \ge 2$---are the principal methodological gaps.

\subsection{Case Study 4: Mapper-based causal abstraction with
calibrated uncertainty}\label{sec:cs_mapper}

\paragraph{Setting.}
Many causal questions in medicine and the social sciences concern
\emph{regimes}: disease subtypes, treatment-responder strata, market
phases. Regimes are typically not pre-specified by covariates but emerge
from the geometry of the data. The Mapper algorithm
\citep{singh2007mapper,carriere2018statistical,dlotko2019ball,
madukpe2026mapper} produces a topological abstraction of a
high-dimensional data set into a simplicial complex whose nodes are
clusters of level sets of a chosen filter $f$; we propose that Mapper
abstractions are the natural carrier of a regime-conditional causal
estimand, and that \citet{lecci2014statistical}'s bootstrap-confidence
programme delivers the calibrated uncertainty that regulated industries
require.

\paragraph{TCDA instantiation.}
(i)~$\Mc$ = a high-dimensional feature space of patient (or asset)
descriptors. (ii)~$G$ = a coarse-grained DAG over Mapper nodes
representing regimes and edges representing admissible transitions or
treatment effects between regimes; we assume $G$ is acyclic for
identification, while permitting cycles inside the data Mapper graph
itself. (iii)~$\Tc = \Map_f$ with a filter $f$ chosen to align with the
causal mechanism (disease severity, time-to-event, an exposure score).
(iv)~$\Cc$ = a stratum-conditional ATE, defined for each Mapper node
$\alpha$ by
\begin{equation}
  \CATE(\alpha) \;=\; \E\bigl[\,Y^1 - Y^0 \;\big|\; X \in \alpha\,\bigr],
  \label{eq:cate_mapper}
\end{equation}
together with its uncertainty.

\paragraph{Bootstrap confidence sets for the abstraction.}
The statistical-TDA programme of \citet{lecci2014statistical} supplies
the inferential scaffolding for $\Tc$ via the bottleneck bootstrap of
\citet{fasy2014confidence,chazal2013bootstrap}: a confidence band
$\mathcal{B}_n^{1-\alpha}$ around the empirical persistence diagram of
the Mapper complex covers the population diagram with probability at
least $1 - \alpha$. Translating $\mathcal{B}_n^{1-\alpha}$ into
uncertainty for \cref{eq:cate_mapper} proceeds in two stages. First,
we bootstrap the Mapper graph to obtain a posterior over regime
boundaries, propagating the cluster-membership uncertainty to each
patient or asset. Second, we use the resulting node-level posterior
weights as inputs to a regime-conditional AIPW estimator. The
combination delivers, for each Mapper node $\alpha$, a point estimate
$\widehat\CATE(\alpha)$ and a confidence interval that accounts for
both the topological (regime-assignment) and statistical
(treatment-effect) sources of uncertainty.

\paragraph{Connection to stratified-space causal heterogeneity.}
The Mapper output is a stratified topological abstraction: nodes are
$0$-strata, edges are $1$-strata, and higher simplices encode
higher-order regime overlaps. The theory of stratified spaces
\citep{hughes2004stratified,banagl2007topological,chazal2013persistence}
is exactly the right language for causal effects that vary across such
strata---sub-domain C3 of \Cref{sec:framework}. The bootstrap pipeline
sketched above is therefore a concrete instantiation of Domains A2,
B1, and C3 simultaneously.

\paragraph{Why this case study matters and what remains open.}
The case study addresses a critique frequently levelled at
Mapper-based causal claims in healthcare and finance: the abstraction
is descriptive, not inferential, and conveys no notion of statistical
confidence. \citet{lecci2014statistical}'s programme dissolves the
critique, provided the abstraction is treated as a TCDA functor. The
open methodological question is the design of filter functions $f$
that are simultaneously (i)~causally informative---i.e.\ aligned with
the mechanism of interest---and (ii)~statistically well-behaved under
the bootstrap. We list this as a research avenue in \Cref{sec:research}.

\subsection{Case Study 5: Spatiotemporal causal attribution via
cubical persistent homology}\label{sec:cs_climate}

\paragraph{Setting.}
Climate-event attribution asks whether a specific extreme event would
have occurred in a counterfactual climate without anthropogenic
forcing. Conventional attribution operates on scalar summaries of the
event (return period, intensity) and so is blind to the spatial
structure of the extreme. Cubical persistent homology of the climate
field captures that structure---the number of connected hot regions,
the topology of pressure ridges, the persistence of low-pressure
cells---and so naturally supports a topological notion of attribution
\citep{strommen2025topological,smith2025multiscale,smith2025harnessing}.

\paragraph{TCDA instantiation.}
(i)~$\Mc$ = the space of gridded scalar fields on $\R^2 \times [0,T]$
(daily precipitation, surface temperature, geopotential height), with
the $L^\infty$ pixel-wise distance. (ii)~$G$ = a coarse climate causal
graph linking forcing (anthropogenic emissions, sea-surface
temperature anomalies, modes of variability) to the spatial event
field; the binary attribution treatment is ``factual climate'' versus
``counterfactual climate'' produced by a model ensemble.
(iii)~$\Tc$ = cubical $\PHk{k}$ for $k \in \{0,1\}$ of the sublevel
sets, with optional silhouette smoothing. (iv)~$\Cc$ = a
counterfactual contrast between persistence diagrams,
\begin{equation}
  \Cc_{\mathrm{attr}}
  \;=\; W_1\!\bigl(\Dgmk{k}(P^{\text{fact}}),
                   \Dgmk{k}(P^{\text{cf}})\bigr),
  \label{eq:climate_contrast}
\end{equation}
or its silhouette functional analogue analogous to
\cref{eq:tate_silhouette}.

\paragraph{The descriptive precedent and the causal upgrade.}
The cosmology literature \citep{sousbie2011persistent,
vandeweygaert2011alpha,cisewski2014nonparametric} provides the mature
descriptive template: persistence diagrams of cosmological density
fields characterise the cosmic web, and bootstrap confidence sets
on those diagrams supply calibrated descriptive uncertainty. The TCDA
upgrade is straightforward in principle: replace cosmology's
``observed universe vs.\ simulated universe'' contrast by climate's
``factual climate vs.\ counterfactual climate'' contrast, and apply
the AIPW machinery of \citet{kim2026topological} treating each ensemble
member as a sample from the relevant interventional distribution. The
stability transfer of \Cref{prop:stability-transfer} controls the
sensitivity of \cref{eq:climate_contrast} to grid resolution and to
the small Hausdorff-scale perturbations that arise from regridding,
bias correction, and interpolation.

\paragraph{Why this case study matters and what remains open.}
The case study realises sub-domains D2 (spatiotemporal causality) and
B1 (stability of causal estimates) of \Cref{sec:framework} on a
high-impact application. The principal remaining gap is identifiability
of \cref{eq:climate_contrast} under unobserved natural variability:
a coupled-model ensemble approximates the counterfactual distribution
but does not eliminate the latent low-frequency oscillations that act
as confounders. \Cref{prop:confounding_signature} suggests a built-in
diagnostic: a persistent $1$-cycle in the residual climate field
\emph{after} adjustment for the assumed forcing covariates would
flag that some mode of variability has not been accounted for, and so
the attribution estimate carries the residual topology of the
unidentified causal structure. Climate TCDA is therefore both a
showcase application and a stress test of the framework: its
methodological success will depend on whether the diagnostics of
sub-domain B2 can be deployed in concert with the estimators of
sub-domains A3 and B1.

\subsection{Synthesis: case studies and the taxonomy}\label{sec:cs_synthesis}

\Cref{tab:cs_summary} collects the case studies of this section, the
seed and bridge works that anchor each, and the taxonomy nodes of
\Cref{fig:taxonomy} that the case study populates.

\begin{table}[t]
\centering
\caption{Case studies and their place in the TCDA taxonomy of
\Cref{fig:taxonomy}. Each case study instantiates the TCDA tuple
$(\Mc, G, \Tc, \Cc)$ of \Cref{def:tcda} in a specific applied or
theoretical setting, engages one or more of the four seed works of
the field, and populates one or more sub-domains of the taxonomy.}
\label{tab:cs_summary}
\begin{tabular}{@{}lp{0.36\linewidth}p{0.28\linewidth}p{0.16\linewidth}@{}}
\toprule
\S & Case study & Seed and bridge works & Taxonomy node(s) \\
\midrule
\ref{sec:cs_tate}        & Topological ATE on biomedical imaging
                         & \citet{kim2026topological,
                                  lecci2014statistical,
                                  soares2020sars,axelrod2022geom}
                         & A3, B1, D3 \\
\ref{sec:cs_cycles}      & Cycle-based discovery in dependence networks
                         & \citet{elyaagoubi2023statistical,
                                  elyaagoubi2023topological,
                                  bando2022causal}
                         & A1, C1, D1 \\
\ref{sec:cs_confounding} & Confounding diagnostics via residual PH
                         & \citet{ibeling2021topological,
                                  mahadevan2025sheaf,
                                  fasy2014confidence}
                         & B2, B3 \\
\ref{sec:cs_mapper}      & Mapper abstractions with calibrated uncertainty
                         & \citet{lecci2014statistical,
                                  singh2007mapper,
                                  carriere2018statistical}
                         & A2, B1, C3 \\
\ref{sec:cs_climate}     & Spatiotemporal climate attribution
                         & \citet{kim2026topological,
                                  sousbie2011persistent,
                                  strommen2025topological}
                         & D2, B1, B2 \\
\bottomrule
\end{tabular}
\end{table}

Three observations emerge from the synthesis. First, every cell of the
TCDA tuple $(\Mc, G, \Tc, \Cc)$ is exercised non-trivially in at least
one case study: $\Mc$ ranges over image spaces, time-series spaces,
residual spaces, feature spaces, and climate-field spaces; $G$ ranges
from simple treatment-to-outcome DAGs through dependence networks to
coarse-grained regime graphs; $\Tc$ ranges over $\PHk{k}$ in
Vietoris--Rips, cubical, dependence-network, and Mapper forms; and
$\Cc$ ranges over Wasserstein contrasts, silhouette ATEs, hypothesis
tests, and stratum-conditional CATEs. The expressive scope of
\Cref{def:tcda} is therefore borne out empirically.

Second, the four seed works each anchor at least one case study, but no
single seed work anchors more than two: the field is not the work of any
one author or group, and a serious TCDA programme must draw on all four
threads in parallel. The relative position of the seed works in the
case-study landscape---\citet{kim2026topological} as the estimator
champion, \citet{ibeling2021topological} as the no-free-lunch theorist,
\citet{lecci2014statistical} as the statistical infrastructure, and
\citet{elyaagoubi2023statistical} as the benchmark architect---is
exactly the four-cornered division of labour anticipated in
\Cref{sec:framework}.

Third, every case study identifies a concrete remaining gap that we
take up explicitly in \Cref{sec:research}: scaling, transportability,
identifiability, filter design, and identification under unobserved
variability. The case studies are therefore not endpoints but launch
pads, and the agenda of \Cref{sec:research} is calibrated against them.
